\documentclass[12pt,a4paper]{article}

% --- Packages ---
\usepackage[utf8]{inputenc}
\usepackage[T1]{fontenc}
\usepackage[english]{babel}

\usepackage{amsmath}
\usepackage{graphicx}
\usepackage{booktabs}
\usepackage{array}
\usepackage{enumitem}
\usepackage{adjustbox}
\usepackage{pgfplots}
\pgfplotsset{compat=1.18}
\usepackage[protrusion=true,expansion=false]{microtype}
\setlength{\emergencystretch}{3em}
\usepackage{url}
\usepackage{hyperref}
\hypersetup{
    colorlinks=true, linkcolor=blue, filecolor=magenta, urlcolor=cyan,
    pdftitle={A Comparability Protocol for Benchmarking Relational Database Access Layers in Express.js},
}
\usepackage{csquotes}

% --- Bibliography (biblatex + biber) ---
\usepackage[backend=biber, style=ieee, sorting=nyt]{biblatex}
\addbibresource{references.bib}

% --- Document metadata ---
\title{A Comparability Protocol for Benchmarking Relational Database\\
  Access Layers in Express.js}
\author{Mateusz Miotk\\
  Faculty of Mathematics, Physics and Informatics,\\
  University of Gdańsk\\
  \texttt{mateusz.miotk@ug.edu.pl}}
\date{\today}

\begin{document}
\maketitle

\begin{abstract}
\noindent\textbf{Context:} Relational access-layer benchmarks can compare unequal
results or report latency at unequal proximity to saturation.
\textbf{Objectives:} We develop and evaluate a protocol for deciding when
configured implementations are comparable and what conclusions their measurements
support.
\textbf{Methods:} The protocol requires task-derived read checks,
cross-implementation equivalence, post-write state validation, a predeclared
configuration-selection policy, and separate capacity and latency operating
points. An open Express.js case study compares nine policy-selected configurations
and two tuned native references on PostgreSQL, MySQL, and five access patterns in a
25-run randomized-block campaign. A second same-host campaign tests selected rank
reversals; a source-located audit applies the protocol to nine external benchmarks.
\textbf{Results:} The largest observed spread is the tested deep fetch:
$7.01\times$ on PostgreSQL and $5.12\times$ on MySQL. With common SQL executed
through raw facilities, the spreads are $1.71\times$ and $2.03\times$; this
compound contrast isolates no single mechanism. Deep-fetch p99 contracts from
18--108 to 2.3--4.9~ms on PostgreSQL and from 28--121 to 1.9--5.5~ms on MySQL at
the stable 50\% capacity target. Four material cross-stack rank reversals recur in
both same-host campaigns; close MySQL insert ranks do not.
\textbf{Conclusion:} The protocol provides an auditable discipline for
establishing and reporting comparability, rather than a claim of universal
rankings or library-intrinsic overhead. Its external audit has one rater and the
adapters lack independent human review, so independent applicability remains to
be tested.
\end{abstract}

\vspace{0.5em}
\noindent\textbf{Keywords:} benchmarking methodology; database access layer; object-relational mapping;
Node.js; response-time tail (p99); empirical software engineering

% \tabledir is set here so the shared results section can locate the generated
% tables; the IST build (paper/ist/) sets it to ../tables/ instead.
\newcommand{\tabledir}{tables/}

\section{Introduction}
\label{sec:introduction}

Express.js is among the most widely used web frameworks for Node.js, and almost every non-trivial
service built on it reaches a relational database through some \emph{access layer}. Access layers
trade raw performance for ergonomics, type safety, and protection from pitfalls such as the N+1 query
problem~\cite{chen2014orm} (Section~\ref{sec:related}). Native drivers expose the wire protocol
directly (\texttt{pg}, \texttt{mysql2}), query builders assemble SQL programmatically (Knex), and
object-relational mappers (ORMs) map rows to objects following patterns such as Data Mapper and
Active Record~\cite{fowler2002patterns}. Node.js back-end performance under load is a recognized
empirical topic~\cite{chaniotis2015nodejs,kuffel2025nodejs,carmo2024backend}, but that literature
varies the runtime and framework, not the access layer.

The most visible comparisons are vendor-authored, and the peer-reviewed ones surveyed in
Section~\ref{sec:related} each leave part of the space uncovered: breadth of products, both engines,
a reported tail, or measurement through an HTTP framework (Supplement Table~S51).
These gaps motivate the coverage of this experiment, not a priority claim.
A more basic gap is methodological: the reported numbers are not established to be
\emph{comparable}, since a layer returning fewer fields or erroring faster than it works can look
faster, a library's chosen query strategy is conflated with its own cost, and a single, arbitrary
load point conflates capacity, tail latency under demand, and per-request latency.
Establishing that comparability is the problem this paper takes up, and it is a
software-engineering one: where such benchmarks inform a technology choice, a comparison not
established to be comparable can misdirect that choice rather than merely misreport a number.

This paper addresses these gaps with a \emph{comparability protocol} (a structured, executable
proposal rather than a validated methodology) and one non-vendor-authored empirical instantiation of
it: an \emph{identical} Express application served through all nine access layers over five selected
access patterns, which stress distinct aspects of the access layer and are not a representative
sample of Express and database workloads. Seven layers run against \emph{both} engines under an
identical schema, dataset, and pool size; two tuned native-driver baselines are disclosed reference
points, not an optimization ceiling. Every (layer, engine, pattern) reports throughput and the
high-load p99 under equal closed-loop demand at a fixed 50-connection point over 25 runs. That point
sits near saturation, where the p99 largely tracks capacity and queueing, so we also report the tail
at \emph{matched utilization}.
The study separates two estimands: the \emph{policy-selected documented configuration}
performance of each layer (RQ1--RQ3), and a \emph{common-SQL raw-path sensitivity} contrast reporting
the residual under identical SQL, which changes query formulation, protocol, preparation, round trips,
and mapping together rather than isolating any one. The full harness (adapters, seed data, an
autocannon-driven~\cite{autocannon} load runner, a specification-derived expected-result oracle,
differential-equivalence checks, and write-state validation) is released as an open replication
package.\footnote{Harness, deterministic seed, all adapters, raw per-cell measurements, and the
table/figure-generating scripts: \texttt{https://github.com/mmiotk/express-db-access-performance}.
The archived release and its DOIs are given in the Data availability statement.}

The experiment answers three research questions, each scoped to the benchmarked
implementations under the specified operating point:
\begin{enumerate}[label=\textbf{RQ\arabic*},leftmargin=*]
  \item \emph{(performance difference):} How large is the throughput and response-time-p99
        \emph{difference} relative to the native-driver baseline, and how does it vary across the
        five access patterns? We keep distinct three quantities a fixed operating point conflates:
        \emph{capacity}, \emph{high-load latency}, and \emph{latency at matched utilization}.
  \item \emph{(backend-stack transfer):} Does the observed layer ordering transfer between the
        tested PostgreSQL and MySQL backend stacks, where changing the DBMS can also change its
        supported adapter and lower-level driver? The question concerns transfer, not causation: the
        design cannot attribute a failure to transfer to the DBMS rather than to the adapter or
        driver the backend brings with it. We answer it at two levels, whether a specific pair
        reverses durably under a stated rule and whether the whole ordering agrees, the latter
        reported descriptively.
  \item \emph{(pattern spread):} Which of the five tested patterns shows the
        largest native-relative spread in this configuration?
\end{enumerate}
Supplement Table~S47 states, per question, the quantity the design estimates and the quantity it
does \emph{not} identify.

This paper's primary contribution is a \textbf{comparability protocol for access-layer
benchmarking}, specified normatively in Section~\ref{sec:protocol} (inputs, stages, cell-admission
criteria, interpretation, and limits).
It rests on three constructs, defined there and named here.
\emph{Semantic admission} requires evidence that a treatment performed the declared task before any
of its timings may be used, so a layer that returns less, writes wrong, or errors faster than it
works cannot appear faster.
\emph{Treatment definition} fixes each treatment by a predeclared, reproducible selection rule,
which prevents result-dependent tuning but establishes nothing about whether the configuration is
optimal, common, or preferred in production.
\emph{Operating-point separation} reports capacity, high-load tail under equal demand, and tail at
matched utilization as \emph{distinct} quantities.
These constituent controls are established performance-evaluation and testing
ideas~\cite{raasveldt2018fair,fruth2022tail,georges2007rigorous}, and CYNTHIA already applies
differential equivalence across ORM implementations~\cite{sotiropoulos2021cynthia}; the contribution
is their integration and operationalization for relational access-layer comparison.
We \emph{exercise} that discipline through an in-experiment ablation and a single-rater audit of
nine external benchmarks; neither is independent validation, and Section~\ref{sec:discussion} sets
out what the exercise does and does not establish.

Two contributions should be kept apart. The \emph{methodological} contribution is the protocol
and its checklist: proposals for how access-layer comparisons should be conducted, reported,
and audited. The \emph{empirical} contribution is one instantiation of them, scoped to the
conditions and software versions of Section~\ref{sec:methodology}. No result is offered as a
property of a library, and the study measures no non-performance quality such as maintainability,
type safety, migration tooling, or ecosystem maturity.

\section{Related Work}
\label{sec:related}

We survey prior work along four \emph{coverage} axes (taxonomy, engine, joint throughput and tail
reporting, and vendor involvement) spanned by this experiment (Supplement Table~S51), and separately
audit which comparability controls each source reports. A structured scoping search, not a systematic
review, of the ACM Digital Library, IEEE Xplore, Scopus, and Google Scholar (2006--July~2026)
retained empirical comparisons of relational access layers or engines, with details archived in the
replication package; being scoping rather than exhaustive, we scope every gap claim to the sources
searched.

\subsection{Object-relational mapping and the access-layer spectrum}
Access layers bridge the object-relational impedance mismatch between the relational model and an
application's object graphs~\cite{ireland2009impedance}, spanning native drivers, query builders, and
ORMs following Data Mapper / Active Record patterns~\cite{fowler2002patterns}
(Section~\ref{sec:introduction}). Torres et al.~\cite{torres2017orm} survey two decades of ORM
patterns; their trade-off (productivity against a hard-to-predict runtime cost) motivates what
we measure while holding the store family (relational rather than non-relational) fixed. The treatment is the configured
access-layer implementation bundle, including its selected API, SQL, protocol, and mapping path;
the study does not isolate a library binary from those choices.

\subsection{ORM performance and data-access anti-patterns}
A parallel line studies ORM performance in reverse, repairing the inefficient SQL frameworks emit.
The costliest is the N+1 query problem (Section~\ref{sec:introduction}); it and related
redundant-query anti-patterns pervade real-world web
applications~\cite{yang2018dbapps}, their impact quantified in ORM-backed
systems~\cite{chen2016redundant,chen2014orm} and often \emph{removable} by query synthesis,
batching, or caching. Closest to our
design, Lorenz et al.~\cite{lorenz2017orm} show the best mapping strategy depends on database
technology, van Zyl et al.~\cite{vanzyl2009optimisations} that it is tuning-sensitive, and a recent
study adds an energy dimension~\cite{bonvoisin2024orm}. Overhead is thus configurable, not fixed, 
as our per-pattern results confirm, but this line never compares drivers, query builders, and ORMs
head-to-head under load, our aim.

Cross-implementation ORM equivalence checking has a direct precedent. CYNTHIA generates an abstract
query, translates it to equivalent queries for five ORMs, executes them over four relational DBMSs,
and differentially compares the results to find correctness bugs~\cite{sotiropoulos2021cynthia};
matching results alone do not establish correctness against an independent task specification.
CYNTHIA is a testing method, not a performance-benchmark admission protocol: it does not use
equivalence to decide which timing cells may enter a benchmark, define a premeasurement rule for
which implementation represents each library, or separate equal-demand capacity effects from
matched-utilization latency. Our claim is restricted to using specification evidence and
cross-implementation equivalence as benchmark admission, joined with predeclared treatment selection
and operating-point separation.

\subsection{Peer-reviewed access-layer comparisons and adjacent lines}
Peer-reviewed head-to-head access-layer measurement is scarce, and each study leaves part of the
space uncovered. Colley et al.~\cite{colley2018orm} illustrate Entity Framework-generated against
hand-written SQL rather than comparing layers; two recent studies compare Prisma with optimized raw
SQL, and Prisma with TypeORM and Sequelize, both only on
PostgreSQL~\cite{prisma_rawsql_2025,orm_compare_jcct_2025}; and a third times internal query stages
of three ORMs rather than client-observed requests~\cite{zhadko2025orm}. Pratama and
Raharja~\cite{pratama2023nodebench} do vary five access products, but framework, access path, and
environment change across cells too, so no access-layer estimand is isolated. A further Express
study optimizes one stack rather than comparing layers~\cite{imamsakti2025async}.

A second, disjoint line benchmarks the layers below and above the access layer without varying it.
Salunke and Ouda~\cite{salunke2024pgmysql} compare PostgreSQL and MySQL under standard OLTP
workloads, the only dual-engine precedent found, but benchmark the \emph{engines themselves}. On the
Node.js side, work establishes Node's web viability~\cite{chaniotis2015nodejs}, tracks back-end drift
under sustained load~\cite{kuffel2025nodejs}, and compares frameworks
broadly~\cite{carmo2024backend}, in each case holding the access layer fixed.

\subsection{Standard benchmarks and benchmarking methodology}
Standard benchmarks define how systems are measured: YCSB~\cite{cooper2010ycsb} for cloud-serving
stores and TPC-C/TPC-E for OLTP. These target the \emph{engine}, not the driver, query builder, or ORM
this study varies, but set the bar for controlled, repeatable load. A methodological literature
explains why single-metric reporting fails: Fruth et al.~\cite{fruth2022tail} and Raasveldt et
al.~\cite{raasveldt2018fair} catalog benchmarking pitfalls, Dean and Barroso~\cite{dean2013tail} show
the tail drives user-perceived latency under fan-out, and Li et al.~\cite{li2014tales} trace its
sources. Such measurement is often ad hoc~\cite{leitner2017exploratory} and hard to reproduce, motivating a harness built for relevance, repeatability, and fairness. Following that guidance, we report
throughput \emph{and} p99 for the same run: a pairing the general literature recommends, which we did
not find in the access-layer benchmarks surveyed here.

\subsection{Practitioner and vendor benchmarks}
A larger body comes from practitioners and vendors, each authored by a party with a stake in a
compared product: Drizzle's Northwind micro-suite and k6 load test~\cite{drizzle_benchmarks},
Prisma's median-latency comparison reporting neither throughput nor any
percentile~\cite{prisma_benchmarks}, and IMDBench over three CRUD operations~\cite{imdbench}.
Together these span more products than the peer-reviewed literature, but none combines
dual-backend-stack coverage, response-time p99, capacity characterization, and frozen
source-located treatment provenance.

\subsection{Positioning}
\paragraph{Audit corpus.} The nine audited sources are those retained by the same scoping search,
under stated criteria rather than convenience: empirical performance comparisons of relational
access layers, or direct PostgreSQL-versus-MySQL comparisons, bearing on the four positioning
properties. Peer-reviewed studies, vendor suites, and practitioner benchmarks were all eligible,
since the audit concerns what a benchmark reports, not who published it. Coding was performed once,
by the author, from a stable locator for every judgement; \texttt{not\_reported} records what the
inspected source does not state, never a finding that the authors omitted the practice. The corpus
is not exhaustive, so every gap statement below is scoped to it: we do not claim that no other
benchmark implements these controls. Eligibility and exclusion criteria, the codebook, the
screening record, and per-judgement evidence are in the online supplement and the replication
package.

In the sources searched, no access-layer benchmark reports all seven protocol dimensions,
unsurprisingly, since the seven are this paper's own construction: the audit measures distance from a
new standard, not failure to meet a known one. The source-located audit records partial controls
rather than treating non-reporting as proof of absence (Supplement Table~S37), and we found no study
combining broad product coverage, both engines, joint throughput and tail reporting, and non-vendor
authorship, each covering at most two or three (Supplement Table~S51): the strongest claim the
scoping search supports, not one of priority. 

But coverage is not the contribution: the protocol operationalizes established experimental principles for
this genre, not a new one. Kounev et al.'s systems-benchmarking
criteria~\cite{kounev2020benchmarking}, Hasselbring's software-engineering
standard~\cite{hasselbring2021benchmarking}, and Raasveldt and M\"uhleisen's fair-benchmarking
guidance~\cite{raasveldt2018fair} supply its principles at a deliberately general level. They leave an
access-layer study to operationalize \emph{what code represents a library}, how competing
implementations are checked for semantic equivalence, and how strategy-sensitive results are
contrasted. Our protocol supplies one concrete operationalization: an arbitrary but predeclared treatment policy; a common-SQL raw-path sensitivity contrast; and a per-cell gate requiring specification-conformant, mutually equivalent responses and verified write state \emph{before} a timing is admitted. Neither element is new in itself, correctness-before-timing having precedents in SPEC validation, TPC ACID, and SQLancer~\cite{rigger2020sqlancer}, and equivalence checking in CYNTHIA above; the contribution is their composition into a discipline for establishing and reporting comparability among competing implementations of one task, stated in Section~\ref{sec:introduction}. Its evidence is the in-experiment ablation and a single-rater retrospective audit of nine external benchmarks (Supplement Table~S37); this demonstrates utility and applicability, not independent inter-rater validation or invention of the underlying principles.


\section{Study Design}
\label{sec:methodology}

In Goal--Question--Metric terms, the study analyzes configured relational-access-layer implementations to compare throughput and response-time p99 under PostgreSQL and MySQL back ends~\cite{basili1988tame}, from the viewpoint of the benchmark analyst evaluating policy-defined treatments rather than a model of observed developer behavior.
This is a \emph{controlled benchmarking experiment} on one synthetic service, dataset, and environment, not inference about ORMs as a population: treatments are assigned and executed automatically rather than observed in situ, factors are fixed by design, run order is randomized in blocks, and each cell is measured repeatedly~\cite{wohlin2012experimentation,kitchenham2002guidelines}.
That design licenses the paired within-campaign analysis and bounds external validity to the single instantiation measured here.
The experiment instantiates the protocol with three decisions: semantic admission before a cell's timing is used, a predeclared treatment-selection policy, and explicit operating points.
Admission gates the \emph{use} of a timing, not always its execution: several admission checks were added after the campaigns they screen (Supplement Table~S48).

The primary RQ1--RQ3 treatment is the \emph{policy-selected documented configuration}: the path selected mechanically from frozen official documentation under the rule below.
The protocol does not privilege documentation; another study may preregister an expert-tuned, usage-frequency, or vendor-recommended policy.
A \emph{common-SQL raw-path sensitivity contrast} then has every portable layer execute the same SQL through its raw facility.
SQL, API level, protocol, preparation, round trips, and mapping still change together, so the contrast measures sensitivity to standardizing the whole bundle; it decomposes nothing and bounds nothing (Supplement Table~S42).

A fixed concurrency imposes equal \emph{demand} on layers of unequal throughput and so drives them to unequal \emph{utilization}.
RQ1 is therefore characterized under three operating conditions, the protocol's \emph{operating-point separation}.
\emph{Equal external demand} is the fixed 50-connection \emph{high-load} point covering the full matrix; a per-pattern concurrency ladder places every pattern's throughput knee at or below 50 connections, so this point sits at high utilization of each pattern's own capacity (median 98--100\%), a capacity-and-queueing regime rather than a sub-saturation latency comparison (Supplement Table~S35).
The two utilization-dependent conditions need a per-layer capacity target and are shown on the deep fetch: \emph{equal utilization} offers each layer a fixed fraction (50/70/85/95\%) of its \emph{own} estimated capacity under the coordinated-omission-corrected open-loop sweep, and \emph{equal compute budget} is the exploratory equal-CPU check.
These answer different questions and need not agree, so all three are reported.
The capacity proxy $\widehat C_{50}$ is the median throughput of the 25-run, 50-connection primary cell and remains an estimate (Supplement Table~S45); the 50\% and most 70\% cells stay below saturation, while 85\% and 95\% are read only as near-saturation trends (Supplement Tables~S21--S22).

\subsection{The comparability protocol}
\label{sec:protocol}
We state the protocol independently of this experiment, then instantiate it.

\paragraph{Comparability, operationally.} Configured treatments are \emph{comparable for a stated estimand} when each satisfies the predeclared semantic-admission and treatment-definition criteria that estimand requires, and the differences that remain between them are the differences the estimand is about.
Comparability is therefore a property of a claim, not of a pair of systems, as this study's selected-configuration and common-SQL estimands illustrate.
The definition deliberately does \emph{not} require identical SQL, round-trip counts, internal algorithms, or framework responsibilities; requiring any of these would suppress the design choices an access-layer comparison exists to measure.
What admission establishes is narrower: every timed treatment performed the declared task, returned the declared result, and left the declared state, so a difference in timing is not a difference in work done.
Admission is a floor, not a certificate of equivalence: admitted treatments stay unequal in ways the protocol records rather than removes, and no result below should be read as comparing otherwise-identical systems.

Two stages are mandatory, \emph{M1} semantic admission and \emph{M2} treatment definition, and five are recommended, \emph{R3} to \emph{R7}; Figure~\ref{fig:protocol} states the inputs, stages, cell-admission rule, and output interpretation.
They are not interchangeable: each addresses a distinct threat to interpretability, and a study may satisfy some and not others, which narrows the admissible claims rather than invalidating the reported numbers.
The stages cover the admission-through-interpretation dimensions addressed here; they are not asserted to be a complete taxonomy of every threat to access-layer benchmarking.
Here M1--M2 cover the full matrix, R4 every pattern, R3/R5/R6 the deep fetch, and R7 the measured source.
The full normative specification (stage definitions, the cell-admission rule, applicability limits of the output comparator, the six-level compliance scheme, and the level each workload cell attains) is in the online supplement, with a machine-readable encoding in \texttt{protocol-checklist.yaml}.

\usetikzlibrary{arrows.meta,positioning,fit,shapes.geometric}
\input{\tabledir fig_protocol}

\paragraph{Protocol development versus protocol validation} The protocol was not specified in full before the measurements it governs.
We separate \emph{development evidence} (the protocol finding defects in this study's own harness) from \emph{internal demonstration} (the worked instantiation reported here) and \emph{external validation}, which has not been performed.
Four elements postdate the accepted campaigns and are therefore retrospective, applied to archived state rather than gating the timed runs (Supplement Table~S48).
That the protocol exposed real defects shows the checks do work, but the defects were found by the author who designed them, so this is not independent validation.

\subsection{Factors and treatments}
We vary three factors.
The \emph{access layer} has eleven levels across three tiers.
Two native drivers (\texttt{pg}, \texttt{mysql2}) and their two engine-specific tuned baselines each run only on their own engine; the query builder Knex and the six ORMs (Drizzle, Prisma, Sequelize, TypeORM, Objection, MikroORM) run on both.
The tuned baselines are disclosed reference points, not selected treatments, leaving \emph{nine} policy-selected layers and $4 + 7\times 2 = 18$ layer$\times$engine cells across 5 patterns, so 90 measured configurations.
We classify a library by its dominant interface, not its self-description (Section~\ref{sec:introduction}).
The nine are a deliberately broad product cross-section, not balanced sampling of the three tiers and not an exhaustive catalogue; inclusion and exclusion criteria are in \texttt{METHODOLOGY.md}.
\emph{Non-vendor authorship} means only that no evaluated library's maintainers were involved; the consequential author choices (selection, schema, pool size, tuning) remain and are disclosed.
All adapter code and treatment assignments were produced by the single author, a limitation analysed in Section~\ref{sec:threats}.

The \emph{backend-stack} factor has two nominal DBMS levels (PostgreSQL~18.4 and MySQL~9.7.1), but switching it can also switch the supported adapter, driver, wire protocol, and defaults, so RQ2 estimates transfer between stacks, not a pure DBMS effect.
The \emph{access pattern} has five levels, realized as HTTP endpoints (Supplement Table~S39).
Supplement Table~S29 records which factors are held constant, vary, or are uncontrolled on this single virtualized host, the reason we report this configuration rather than general library performance.

\subsection{Objects: schema, workload, and data}
The workload is a minimal blog domain (\texttt{authors} $1\!-\!*$ \texttt{posts} $1\!-\!*$ \texttt{comments} $*\!-\!1$ \texttt{authors}).
The five endpoints implement five selected access patterns~\cite{cooper2010ycsb}: a keyset page (primary-key cursor, not a large \texttt{OFFSET}), a single-row read, a one-to-many list, the deep fetch (a post with its author and every comment, each with its own comment-author), and a per-author aggregation (Supplement Table~S39).
A deterministic seed (2{,}000 authors, 100{,}000 posts, 1{,}000{,}000 comments) is loaded by the native driver, independently of any layer under test, so both engines receive byte-identical row values.
The working set fits in memory, so measurements emphasize the access layer's per-request instruction, protocol, and mapping cost, while engine-side execution, buffering, locking, and logging remain in every measurement~\cite{harizopoulos2008oltp}.
Each request's target identifier is drawn uniformly from the seeded range, spreading load across the working set; a skewed, production-like distribution is a planned variant (Section~\ref{sec:threats}).

\subsection{The adapter contract and N+1 control}
\label{sec:adapter}
Every access layer implements the same factory interface, so the server and load runner are layer-agnostic and each returns an \emph{identical result shape}.
Each avoids the deep-fetch N+1 problem by its selected mechanism, a hand-written join for the native drivers, Knex, and Drizzle, and relation eager loading for the data-mapper ORMs.
The comparison is thus of each layer's \emph{selected configuration} under a fixed adapter-selection rule, not a fixed query plan: a measured difference reflects the SQL generated, the round-trip count, and result materialization, all captured by server-side logging (Supplement Table~S2).

The API was fixed by a rule recorded before the accepted campaigns (2026-07-22; Supplement Table~S48) and applied unchanged: the eager-loading API each layer's version-compatible official documentation presents first, ties broken toward the lower-level API the library documents as its base layer.
This is an intentionally artificial, predeclared \emph{policy}, not evidence that the chosen API is performance-optimal, the most common production choice, or what an expert would deploy; it is one instance of the protocol's predeclared treatment rule, which privileges none (corroboration and limits: Table~S33, Section~\ref{sec:threats}).
Reproducible selection is also not yet independently \emph{reproduced}: all four result-blind review slots stand pending, so whether a second analyst applying this rule to the same frozen pages would select the same APIs is untested.

Supplement Table~S16 and the rubric (\texttt{notes/\allowbreak documentation-selection.md}) record, per layer, the selected API with its documentation page, version, and URL, its round-trip count, the documented alternative not used, that logging, hooks, and validation are off, and the freeze date.
The artifact ships the exact pre-freeze page captures with timestamps and hashes, so reconstruction does not depend on today's website; a capture proves the page state at its recorded timestamp, not that it was unchanged through the freeze.
The conflict and tie-break rules are in the online supplement.

For the protocol's \emph{common-SQL raw-path sensitivity contrast}, every adapter also executes the same two-statement deep-fetch SQL and row mapping through its raw-SQL facility, with the same \texttt{EXPLAIN} plan per engine; server-side capture confirms the statements are byte-identical across every layer on both engines, while the raw facility, preparation, and wire protocol still differ (Supplement Table~S42).
A secondary loading-strategy sensitivity switches only between documented, byte-identical deep-fetch alternatives; it does not redefine the primary policy-selected treatment or estimate a library-intrinsic effect (Supplement Table~S18).

Read admission has an expected-result layer independent of any timed reference: \texttt{bench/verify-spec.mjs} replays the documented seed PRNG and fan-out fixture in memory, without querying a database or importing an access layer, and derives exact fields, values, ordering, graph membership, and aggregates.
All nine engine-compatible adapters pass 36{,}540 expected-result checks per engine, 73{,}080 in total.
A separate differential layer establishes that implementations agree, and a companion check validates database \emph{state} rather than HTTP status, confirming for every adapter on both engines that the primary insert wrote exactly the requested values and identifier and that the exploratory transactional method commits atomically or rolls back to no partial state; all declared methods pass.
Both layers provide finite evidence for the declared task and tested inputs; neither is presented as proof for arbitrary inputs.
The canonical constructors, per-adapter scope, fan-out and timestamp conventions, and fixture details are in the online supplement and \texttt{experiments/METHODOLOGY.md}.
The read cross-check caught two defects during development (a fan-out aggregation bug and a driver result-shape mismatch; Section~\ref{sec:discussion}).

\subsection{Controls}
Four controls hold the shared conditions fixed (connection pool, logical task, physical dataset state, and observer) so the \emph{configured access-layer implementation bundle} is the only \emph{deliberately varied} factor.
They do not isolate the library alone: the bundled differences in SQL, query count, protocol, loading strategy, and hydration are the treatment, not confounds to remove (Supplement Table~S42).
\emph{(i)~Pool size} is fixed at ten connections for every adapter, so the comparison reflects the access-layer choice, not pool tuning~\cite{gupta2017pooling,laigner2021microservices}; because the generator opens 50 connections against that pool, requests queue throughout, with instrumented layers reporting mean pending depth 29--38, unobserved for Prisma and MikroORM, which expose no pool metric.
\emph{(ii)~Common task semantics}: each endpoint implements the same declared input--output task, while SQL text, query count, plan, protocol, loading strategy, and hydration remain measured properties of the treatment.
Aggregation is the exception, with a common correlated-subquery formulation.
\emph{(iii)~Write isolation}: the declared reset removes generated rows before warm-up, measurement, and the next cell.
A post-review state oracle and a cross-engine gate exposed two defects in the historical harness, an allocator that could let generated rows persist and a shared-fixture PRNG that gave different comment-author assignments across engines; the corrected harness pins the allocator, verifies exact base and fan-out counts by idempotent preflight, and materializes the fixture once so the joined representation is byte-identical across engines.
The affected pilot was rejected and a corrected-state campaign remeasured all five patterns for all 18 compatible treatment--stack pairs under one version of the instrument; the historical MySQL range-scan, aggregation, and insert cells are not used as clean evidence.
Reset mechanics, sentinel values, and fixture hashes are in the online supplement and \texttt{experiments/METHODOLOGY.md}.
\emph{(iv)~Observer separation}: the HTTP load generator runs in a process separate from the server.

\subsection{Measurement procedure}
Load was generated with \texttt{autocannon}~\cite{autocannon}, a closed-loop (constant-concurrency) HTTP/1.1 generator~\cite{schroeder2006open} recording a latency histogram.
The accepted corrected-state primary matrix contains 90 configurations with 25 repeated runs each; the superseded campaign remains provenance only.
Each replicate randomizes adapter--engine blocks and boots one fresh server process per block, so no application process crosses an adapter, engine, or replicate boundary; read endpoints within a block share that process, and database buffers plus host state are not reset.
Each endpoint receives a discarded 15-second warm-up followed by one 12-second measurement at 50 connections, and a forward/reverse-order write-cell check found no detectable order effect.
The campaign records its ID, deterministic order seed, environment fingerprint, and exact source manifest; process, teardown, and rebuild mechanics are in \texttt{experiments/METHODOLOGY.md}.

Within a replicate, every layer and engine draws target identifiers from one generator seeded on endpoint and replicate, so all face the same identifier distribution, though not an identical served sequence; the paired analysis is licensed by the block design, not by sequence identity.
The inferential unit is the endpoint-level run, individual requests being subsamples, and we report the median of the 25 replicate values of throughput and of the run-level percentiles p50, p90, p97.5, and p99.

Throughout, \emph{p99} denotes this \emph{median run-level p99}, not the p99 of the pooled request distribution, which the design does not estimate.
Ten 60~s re-measurements of the deep fetch give 12-versus-60-second rank correlations of 1.00 on both engines (Supplement Table~S23): endpoint- and campaign-specific sensitivity evidence, not a general claim that 12~s suffices for p99 estimation.
We sampled server CPU and peak RSS over the process tree, with database and load-generator CPU separately.
Inserts were measured under each engine's \emph{default} durability, the regime a deployment runs unless relaxed; a labelled secondary run relaxed all of these as a compound durability sensitivity (Supplement Table~S3).
To characterize behavior under load we swept the deep fetch on a separate 1--256-connection ladder for every layer and engine~\cite{yu2014concurrency}, and a constant-arrival generator drove it under a fixed offered rate rather than fixed concurrency, measuring from intended send time, a coordinated-omission-corrected design~\cite{tene2015latency,krishnamachari2026stats} checking the closed-loop tail's sensitivity to a different load model (Section~\ref{sec:results}).

\subsection{Analysis}
We separate \emph{primary} outcomes, which carry the research-question claims, from \emph{secondary} and \emph{exploratory} outcomes reported descriptively as controls and sensitivity analyses (Supplement Table~S34); the primary outcomes are throughput and the closed-loop p99 on the five patterns against both engines at the fixed operating point.
Because absolute throughput is hardware-specific, we analyze \emph{relative} differences within an engine, a pattern's \emph{spread} being the untuned native driver's median throughput divided by that of the slowest layer on that pattern.
We report medians over the 25 runs, the coefficient of variation as a stability signal~\cite{georges2007rigorous,laaber2019microbenchmarking}, and a fixed-seed percentile bootstrap 95\% CI for the median; these intervals capture within-campaign, run-to-run variability on this host and configuration, not variation across machines, versions, or deployments (Section~\ref{sec:threats}).
Because the design is a \emph{randomized block}, adjacently ranked layers are compared with \emph{paired} tests on their per-replicate throughput ratios, a two-sided sign-flip permutation test cross-checked with Wilcoxon signed-rank, reporting the geometric-mean paired ratio, its bootstrap CI, the paired dominance, and a bootstrap-interval equivalence assessment against an author-selected $\pm5\%$ margin for the closest pair (Supplement Table~S56); we foreground these effect sizes and interval estimates over $p$-values~\cite{arcuri2014hitchhiker,cliff1993dominance}.
Because the seven adjacent pairs follow the \emph{observed} ranking, their significances are a post-hoc, descriptive robustness check, not a confirmatory family, and the Wilcoxon cross-check reuses the same data and design rather than confirming it independently.
All tests operate on the 25 paired run-level values; no interval is computed over pooled requests.
Randomization is over \emph{measurement order}, so it supports temporal balance, not causal identification: the treatments are pre-existing implementations, not interventions assigned to units, and no result should be read as a causal effect of a library on throughput.
Rank correlations are likewise descriptive, the seven layers being the benchmarked population of configurations rather than a random sample.
Adjacent-layer p99 gaps of one or two milliseconds sit at the millisecond measurement floor and are not read as real; conclusions rest on the large-ratio patterns, chiefly the deep fetch.
The exchangeability assumption and its argument, the full construction, and the statistical-methods notes are in the online supplement and \texttt{notes/supplement-methods}; paired p99 comparisons are Supplement Table~S30.

\subsection{Experimental setup}
All measurements were taken on a single host, a virtualized Intel Xeon Gold~6140 instance (16~vCPUs, single NUMA node, 126~GB RAM; Linux~6.17; Node.js~24.18.0; Express~5), PostgreSQL~18.4 and MySQL~9.7.1 local.
Fourteen labelled secondary experiments ran 2026-07-16 to 2026-07-18, the accepted corrected-state primary matrix overnight on 2026-07-24, and the same-host validation campaign on 2026-07-25, each dated by its archived environment capture (Supplement Table~S48).
The secondary experiments therefore predate the accepted primary matrix rather than sharing its window, and the clean-evidence scoping applies to them as it does to the superseded pilot.
Each library is pinned to the \emph{latest stable release compatible with the harness adapter contract} as of the freeze date (2026-07-15), recorded from the committed lockfile and an automated environment capture (Supplement Table~S11); only Sequelize invokes the earlier-release exception.
Because access-layer performance is version-sensitive, this pins a dated snapshot rather than a permanent ranking (Section~\ref{sec:results}).

\subsection{AI-assisted research tooling}
\label{sec:ai-tooling}
AI assistance covered research software and data analysis, not language editing alone, so it is documented here as part of the method; the publisher's separate manuscript-preparation declaration appears before the references.
The sole assistant was Claude (Anthropic), in three model versions; no other vendor's assistant contributed software, analysis, or measurements, and OpenAI Codex was used only after the manuscript was otherwise complete, for two independent critical pre-submission reviews.
Assistance covered the eleven adapters and harness scaffolding, the table and figure generators, the statistical estimators, and manuscript prose, and extended beyond implementation to proposing candidate analysis procedures, each a named published method, among which the author selected and verified.
Author review is not the safeguard: each category carries a check independent of the generated code, namely unit-tested estimators, per-cell tracing of every measured value to an archived record, and the specification-derived oracle gating the adapters.
Their limit is exact. They establish that the estimators compute what they claim and that each measured number came from an archived measurement; \emph{no} check bears on whether the chosen procedures suit the question, or whether the claims drawn from them are sound, and those rest on author judgement alone.
Per-commit provenance, the affected components, the individual validation checks, and the full extent of assistance in analysis are in the online supplement.

\subsection{Replication package}
The harness is released so the full matrix re-runs with one command against a containerized or local PostgreSQL and MySQL.
Two reproduction statuses are kept distinct: an author-run \emph{reconstruction} regenerates the archived tables from the archived raw records and verifies their hashes, repeatable by a reader without a database, whereas \emph{independent reproduction} would require another team re-executing the measurements on its own hardware and none has been performed (Supplement Table~S31).

\section{Results}
\label{sec:results}

Tables report throughput (req/s, higher better) and tail latency (p99~ms, lower
better) per access layer and engine. The two consequential patterns are in the
main text (deep fetch, Table~\ref{tab:deep_fetch}; insert,
Table~\ref{tab:write}); the point read, keyset range scan, and aggregation are in
Supplement Tables~S12, S13, and~S15, all from the run of
Section~\ref{sec:methodology}. Each native driver has an untuned policy-selected row
(\texttt{pg}, \texttt{mysql2}) and tuned (\texttt{pg-tuned},
\texttt{mysql2-tuned}), the tuned rows marking a hand-tuned reference point, not one
attained without code changes. We compare \emph{relative}
differences within an engine. Two estimands recur: the \emph{performance of the selected
configured treatment} (RQ1--RQ3), which is not library-only overhead, and the
\emph{common-SQL raw-path sensitivity contrast} (Supplement Table~S14), the residual spread
when every layer runs common SQL through its raw facility (Section~\ref{sec:threats}).

% One combined table per access pattern (throughput + p99), per layer and engine.
% \tabledir is set by each main file so this shared section resolves without TEXINPUTS.
\input{\tabledir deep_fetch}
\input{\tabledir write}

\subsection{RQ1: performance difference relative to the native baseline}
The native implementation leads nine of the ten engine--pattern comparisons:
\texttt{pg} or \texttt{pg-tuned} is fastest on all five PostgreSQL patterns, and
\texttt{mysql2} on four of five MySQL patterns. The exception is the MySQL insert,
where the four fastest treatments span 1{,}766--1{,}816~req/s with broadly overlapping
intervals and the native driver ranks fourth (Table~\ref{tab:write}); that ordering is
practically unresolved rather than a defeat, and we do not read a leader from it. The current-generation ORMs sit mid-to-low on the read patterns; Prisma is near the bottom of point reads and aggregation on both engines (Table~\ref{tab:deep_fetch}; Supplement Tables~S12 and~S15), retiring an earlier version's headline of Prisma tying the native driver at a large CPU premium (Section~\ref{sec:discussion}).

\texttt{pg-tuned} and \texttt{mysql2-tuned} track their untuned counterparts within about
$2\%$ on the deep fetch (\texttt{pg-tuned} 3{,}700 vs.\ \texttt{pg} 3{,}638;
\texttt{mysql2-tuned} 2{,}398 vs.\ \texttt{mysql2} 2{,}416). The one material gain is
PostgreSQL aggregation, where named prepared statements add $11\%$ (7{,}583 vs.\ 6{,}807).
On MySQL the binary-protocol baseline is \emph{slower} than \texttt{mysql2} on every pattern
(by $0.4$--$4.7\%$), so \texttt{execute()} buys nothing at this operating point; both tuned
baselines are therefore disclosed reference points, not optimization ceilings.

On the \textbf{point read} the native-relative spread is $2.83\times$ (PostgreSQL)
and $2.14\times$ (MySQL) (Supplement Table~S12). On the \textbf{deep fetch}
(a post, its author, and all comments with their authors) it is the widest
measured --- $7.01\times$ on PostgreSQL (\texttt{pg} 3{,}638, 95\% CI
[3{,}602--3{,}662], down to MikroORM's 519~[512--528]) and $5.12\times$ on MySQL
(\texttt{mysql2} 2{,}416~[2{,}413--2{,}425] down to MikroORM's 472~[464--475]). The
native driver leads on both engines, the data-mapper ORMs (especially
MikroORM) trailing furthest (Table~\ref{tab:deep_fetch}); the comparison is paired
(Section~\ref{sec:methodology}), every adjacent MySQL step large relative to its
bootstrap interval (Supplement Table~S36). \textbf{Aggregation} is the least
layer-sensitive pattern: the native-relative spread is only $1.82\times$
(PostgreSQL) and $1.87\times$ (MySQL) (Supplement Table~S15). Every layer forwards
a single query for this pattern, which also avoids the deep fetch's nested-graph
materialization; the design does not isolate either feature as the cause of the smaller spread.

\subsection{Common-SQL raw-path sensitivity contrast}
Server-side capture confirms common statement text and database plans across raw paths on both engines; API level, wire protocol, preparation, round trips, and mapping can still differ (Supplement Table~S42). The deep-fetch native-relative spread is then far smaller --- $1.71\times$ (PostgreSQL) and $2.03\times$ (MySQL) against the selected-configuration $7.01\times$ and $5.12\times$ (Supplement Table~S14). Thus the selected paths differ more than these common-SQL raw paths; the compound contrast assigns no share to any mechanism. A
loading-strategy sensitivity check (Supplement Table~S18) argues against a mere
selected-strategy artifact: forcing the join-selected ORMs (Sequelize, MikroORM) to
select-in makes them slower, and Objection's batched-select-in selected path runs about
$1.6\times$ slower than a single join on MySQL. Objection's result is therefore
sensitive to the documented loading option; this comparison does not isolate a library-only cost.

\subsection{RQ2: backend-stack transfer of the ranking}
Switching the backend changes a bundle that can include the DBMS, supported adapter,
lower-level driver, wire protocol, and defaults. RQ2 therefore concerns
\emph{backend-stack transfer}, not a causal DBMS effect. In the primary campaign,
the seven portable layers have stack-ratio heterogeneity on every pattern (blocked
interaction $F=127.1$--$458.5$, permutation $p<5{\times}10^{-5}$, the $1/(B{+}1)$ resolution floor at $B=20{,}000$; Supplement Table~S28),
but this within-campaign test does not establish durable rank reversals.

A second same-host campaign re-runs 25 repetitions of the seven portable layers for deep
fetch, aggregation, and insert under an independently drawn block order (Supplement
Table~S46). It re-randomizes execution order only: the identifier generator is seeded on
endpoint and replicate, so the campaign replays the same draws and cannot detect
sensitivity to the workload sample.
Deep-fetch ranks reproduce exactly on both stacks (Spearman $\rho=1.00$, maximum
median drift $2.3\%$); aggregation reproduces 5/7 PostgreSQL and 7/7 MySQL positions
($\rho=0.96$ and $1.00$). MySQL insert is less stable: only 3/7 positions reproduce,
$\rho=0.82$, the leader changes from Drizzle to Knex, and maximum between-campaign median drift reaches
$13.6\%$. We therefore promote a cross-stack reversal only when it occurs in both
campaigns, exceeds the predeclared $5\%$ practical margin on both stacks, and paired
bootstrap intervals exclude equality in opposite directions in both campaigns.
Four pairs meet all criteria: on deep fetch, Prisma exceeds Objection and Sequelize
on PostgreSQL but trails both on MySQL; on aggregation, Prisma exceeds Objection on
PostgreSQL but trails it on MySQL; on insert, Prisma exceeds MikroORM on PostgreSQL
but trails it on MySQL. All four therefore involve one treatment, and that treatment is
the one whose lower-level driver differs between the stacks: Prisma runs the MariaDB
adapter on MySQL because it ships no \texttt{mysql2} adapter, while every other layer
uses \texttt{mysql2} (Section~\ref{sec:threats}). The recurrence is consistent with
backend-stack sensitivity as RQ2 defines it, but it cannot be separated from that driver
substitution, and it is not evidence about ORMs as a group. Other changes in close ranks,
including the identity of the MySQL insert leader, remain exploratory.

The magnitude of stack transfer is also treatment-specific. Primary-campaign
PostgreSQL$\div$MySQL ratios range from $1.10$ to $1.84\times$ on deep fetch,
$1.27$ to $1.82\times$ on aggregation, and $1.55$ to $3.06\times$ on insert
(Supplement Table~S28). These ratios and the recurring reversals show that the
observed ordering cannot be assumed to transfer unchanged between the two tested
stacks. They do not allocate the difference to the DBMS, adapter, driver, protocol,
or defaults.

The primary MySQL insert medians illustrate why fine ordering is not a headline:
Drizzle, Objection, Knex, and the native driver lie between 1{,}766 and
1{,}816~req/s, while several replicate distributions are visibly multimodal or
heavy-tailed and reach $13.6\%$ within-campaign CV --- numerically equal to the between-campaign drift above, but a different quantity (Supplement Figure~S1). Prisma is
the slowest policy-selected treatment at 894~req/s, giving a $1.98\times$
native-relative spread, but the broad interval on that spread is
[1.65--2.09]. Under relaxed durability the native MySQL and Knex paths gain
$2.42\times$ and $2.07\times$, whereas gains differ by treatment; concurrent
wait-event sums are consistent with a shared commit-path contribution but are not
a causal time fraction (Supplement Tables~S3 and S19). PostgreSQL insert spans
$3.18\times$ (6{,}018 to 1{,}894~req/s) and changes by at most $14\%$ in the
relaxed-durability sensitivity. These compound interventions do not decompose
backend-stack cost.

\subsection{RQ3: which pattern spreads widest}
The primary native-relative spreads on PostgreSQL are deep fetch
($7.01\times$), keyset range scan ($4.22\times$), insert ($3.18\times$), point
read ($2.83\times$), and aggregation ($1.82\times$). On MySQL they are deep fetch
($5.12\times$), range scan ($3.88\times$), point read ($2.14\times$), insert
($1.98\times$), and aggregation ($1.87\times$). Thus the tested deep fetch has the
largest spread and aggregation the smallest on both stacks. A full concurrency
sweep places the median knee at 8--32 connections and the median utilization at 50
connections at 98--100\% for every pattern--stack combination; the lowest
individual values are 87--88\% for two PostgreSQL patterns (Supplement Table~S35).
The fixed point is therefore a high-utilization comparison throughout, although it
is not an identical utilization fraction for every layer. An exploratory fan-out
sweep from 0 to 500 comments finds a roughly fixed multiplier rather than a spread
that grows with breadth; it does not vary graph depth, schema, or query mix.

\subsection{Tail latency}
At equal closed-loop demand, deep-fetch p99 spans 18--108~ms on PostgreSQL and
28--121~ms among the policy-selected layers on MySQL (Table~\ref{tab:deep_fetch}). The run is the analysis unit; a
12~s slow-layer run supplies only about 60 observations to its top percentile, so
we use the median of 25 run-level p99 values and show their raw spread. Ten
60-second repetitions preserve the practical ordering, with 12-versus-60-second rank correlation $1.00$ on both engines; the largest median shift is $5.6\%$ on PostgreSQL and $3.8\%$ on MySQL (Supplement Table~S23). This is
endpoint-specific sensitivity evidence, not a general sufficiency claim for a
12-second p99 window.

At the stable nominal 50\% capacity target, coordinated-omission-corrected p99
contracts to 2.3--4.9~ms on PostgreSQL and 1.9--5.5~ms on MySQL; most 70\% cells
remain similarly small (Supplement Tables~S21--S22 and S43). Re-expressing the
targets against the full-sweep maximum places nominal 50\% at 46.6--51.0\% and
nominal 70\% at 65.2--71.5\%; repeated sweep curves widen those ranges only to
46.2--52.2\% and 64.6--73.0\% (S45). The 85\% and 95\% results are unstable
near-saturation sensitivity observations and do not support a convergence claim.
The stable 50\% result shows that the large equal-demand p99 spread contains a
material capacity-and-queueing component; it does not identify its share or imply
a deployment-general sub-saturation bound.

The two operating points also \emph{order the layers differently}, which is the
sharpest consequence of separating them. Between the equal-demand and 50\%-utilization
p99 orderings, Spearman $\rho$ is $0.64$ on PostgreSQL and $0.29$ on MySQL
(Supplement Table~S43). The native driver moves furthest: fastest on the tail at equal
demand on both engines, it falls to third of eight on PostgreSQL and seventh of eight on
MySQL once each layer is held at half its own capacity. On MySQL the reversal is
unambiguous --- \texttt{knex} 1.9~ms [1.8--2.2] against \texttt{mysql2} 4.7~ms [4.0--5.6],
non-overlapping over five runs --- while on PostgreSQL the native and leading portable
ranges overlap (\texttt{pg} 2.6~[2.2--3.0], \texttt{knex} 2.3~[2.3--2.7]), so only the
MySQL inversion is claimed. Reporting a tail at one arbitrary load point would have
recorded the native driver as the tail leader on both engines. This is direct evidence
that the operating-point separation stage changes a conclusion rather than only its
caveats, and it is why capacity and tail are reported as distinct quantities.

\subsection{Measurement stability and effect sizes}
Deep-fetch throughput CV reaches $4.5\%$, much smaller than its multi-fold spreads.
The largest primary throughput CV is $13.6\%$ on MySQL insert, whose raw replicates
are shown because a median and interval can hide regime mixtures. All 90 cells
contain exactly 25 run-level samples and zero timed errors, timeouts, or non-2xx
responses. One pre-warm-up startup failure was retried once and recorded; no timed
sample preceded that retry. The fail-closed roster validation requires 90$\times$25
observations before accepting the file.

Predeclared native-driver contrasts are foregrounded; ranking-selected adjacent
comparisons remain descriptive. Paired bootstrap intervals characterize only
within-campaign variation, while the same-host validation campaign is reported
separately. Differences of about 1~ms in p99 and throughput ratios inside the
predeclared $\pm5\%$ practical margin are not promoted as substantive ordering.

\section{Discussion}
\label{sec:discussion}

\paragraph{Version-sensitivity}
An earlier campaign found Prisma (then 5.22, on an in-process Rust query engine) tying the native
driver on the deep fetch while drawing about five application cores. The \emph{identical} harness on
the re-frozen toolchain did not reproduce it: application CPU now stays near one core and Prisma sits
at mid-to-low throughput (Section~\ref{sec:results}). The non-reproduction is a fact, its cause a
hypothesis, since the re-freeze moved the whole toolchain at once, so the headline is
version-sensitive~\cite{vanzyl2009optimisations}.

\paragraph{Implications for practice}
RQ1 and RQ3 agree the selected-configuration difference is specific to the access \emph{pattern}:
the native-relative spread is widest on the deep fetch, where a layer must hydrate rows and resolve
associations rather than merely issue SQL (the object-relational impedance
mismatch~\cite{ireland2009impedance} made concrete), and narrowest on aggregation. Round-trip count
alone does not explain throughput: single-query MikroORM is slowest on the PostgreSQL deep fetch
while two-query Drizzle outruns four-query Prisma (Supplement Table~S2;
Table~\ref{tab:deep_fetch}). Result-graph materialization and the other bundled implementation
differences are plausible contributors, but the study does not isolate them.

\paragraph{Single-row inserts and backend-stack transfer}
Single-row insert behavior is both stack- and treatment-specific, and the durability and wait-event
controls reported there are consistent with a shared commit-path contribution, not with attributing
a performance floor or a fixed share to the DBMS (Section~\ref{sec:results}).

The second same-host campaign preserves PostgreSQL insert ranks but few MySQL positions
(Section~\ref{sec:results}), so we claim no durable fine ordering for MySQL insert. The recurring, material insert reversal
is narrower: Prisma exceeds MikroORM on PostgreSQL and trails it on MySQL in both
campaigns, with paired intervals excluding equality in opposite directions
(Supplement Table~S54).

Every promoted reversal involves Prisma, whose MySQL path uses the MariaDB adapter rather than
\texttt{mysql2}, so backend-stack sensitivity here is inseparable from that driver substitution and
describes one treatment, not ORMs as a class. \emph{Observed:} every reversal clearing the bar
involves the one layer whose MySQL driver differs. \emph{Interpretation:} at the stated bar only the
asymmetric treatment reverses durably, while whole-ordering agreement stays imperfect either way,
reported descriptively rather than as a further reversal claim. \emph{Not identified:} the design
cannot attribute the reversals to the DBMS rather than the adapter, the driver, or Prisma's
implementation, because holding out the treatment holds out the substitution; the hold-out followed
observing that all four involved one treatment, and its criterion is a design fact fixed
independently of the outcomes. This supports backend-stack sensitivity in the tested
implementations, leaving cross-host transfer unresolved. A five-replicate, one-layer-per-tier
six-statement transaction remains exploratory and does not extend the single-row-insert finding to
writes (Supplement Table~S8).

\paragraph{Methodological lessons: a checklist for access-layer benchmarking}
Several failures changed or could have changed the comparison; we distil them into a reusable
checklist (\texttt{notes/supplement-methods}): an \textbf{aggregation fan-out} and a driver
result-shape mismatch, both caught by specification and differential admission; \textbf{OFFSET
pagination} collapsing on MySQL, easily misread as an access-layer weakness; \textbf{write-state
contamination}, where MySQL allocated inserts below the cleanup floor; and a \textbf{query-formulation
confound} where a mis-scoped pre-aggregation re-scanned the whole table, a query-plan artifact
masquerading as an access-layer effect. The write path needs its own gate: MikroORM~7's dropped
\texttt{persistAndFlush} made every insert throw a 500 faster than a real one, briefly \emph{leading}
MySQL until the non-2xx counter flagged it. Because a 2xx alone would miss a \emph{silently}
incorrect write, the gate also validates post-write state. None is schema-specific; each extends the
database-benchmarking pitfalls of Fruth et al.~\cite{fruth2022tail} and Raasveldt et
al.~\cite{raasveldt2018fair} to the access-layer sub-genre.

\paragraph{Evidence for the protocol contribution}
Each stage had a concrete consequence: an independent seed specification caught shared state and
fixture defects that mutual equality would miss, treatment provenance made the loading-policy
sensitivity auditable, capacity characterization changed the tail interpretation, and the common-SQL
path narrowed but did not decompose the selected-treatment spread (Supplement Table~S50). Applied to
the nine external sources, the codebook yields claim licenses rather than a binary quality score
(Supplement Table~S37).

This demonstrates that the protocol is \emph{executable}: it admits and rejects concrete cells,
changed the reported interpretation at every stage exercised here, and codes an external corpus to a
reproducible per-stage result.
It does not demonstrate \emph{reusability by another person}, and two limits sharpen that: no source
was coded \emph{satisfied} on any stage, so the codebook's upper range is untested, and agreement
against this rater would measure agreement with the protocol's own designer.
The codebook, audit record, agreement scorer, and result-blind forms are released so the question can
be settled against this study.
The claim is therefore a \emph{candidate} discipline and reusable artifact, demonstrated once: not
invention of the constituent controls, not a complete theory of benchmarking validity, and not a
validated methodology.

\paragraph{Practical guidance}
The synthesis follows from RQ1--RQ3 and describes the benchmarked implementations in the tested
configuration (Section~\ref{sec:threats}), not access-layer categories in general.
Because the fan-out sweep varies breadth (0--500 children) but not depth, schema, or query mix, the
deep-fetch peak locates a difference among five probes on one synthetic schema: an observation to
re-measure on the target workload, not a law.
Where a hot path resembles the tested deep fetch the access layer is consequential and worth
benchmarking, the thin paths (native driver, Knex) leading here for throughput and for p99 at high
load, though at matched utilization the MySQL native driver falls to seventh of eight
(Section~\ref{sec:results}).
Where the tested workload is read-simple, or for the tested MySQL single-row insert, the observed
layer spread is smaller.
These are performance findings only: the choice also turns on qualities this study does not measure
(Section~\ref{sec:introduction}), which may outweigh them, and the common-SQL results are a compound
sensitivity contrast on the selected-configuration cost, not a licence to use raw SQL.

\section{Threats to Validity}
\label{sec:threats}

We organize the threats along the four standard validity
categories~\cite{cook1979quasi,wohlin2012experimentation}.
The supporting measurements and mechanics for each are in the online supplement; every limitation is
stated here.

\paragraph{Construct validity of the measures}
Our dependent variables are HTTP-level throughput (req/s) and tail latency (p99) at the Express
endpoint, not query-execution time at the driver boundary.
This is deliberate, capturing the full request round trip, though overhead attributed to the
\emph{access layer} may partly reflect Express routing and serialization.
We hold those near-constant: every adapter satisfies a documented \emph{adapter contract} and funnels
rows through shared canonical constructors whose standalone cost is too small to explain the
multi-fold gap, and a fixed-payload endpoint puts the shared framework-and-serialization floor at
least three times above the fastest deep-fetch cell, so no cell is framework-bound.
The specification-derived oracle and the differential checks establish finite conformance and mutual
equivalence (Section~\ref{sec:methodology}); the differences arise in how each layer reaches an
identical response, not in what is returned.
The reported single-row-insert finding uses each engine's \emph{default} vendor-standard durability,
not one we relaxed, so those insert spreads (Supplement Table~S52) describe the tested default
deployment rather than a tuning choice; the relaxed-durability sensitivity is a compound
intervention and does not decompose its cause (Section~\ref{sec:discussion}, Supplement Table~S3).

\paragraph{Construct validity of the treatment}
RQ1's treatment is each layer's \emph{policy-selected documented} implementation and eager-loading
strategy: an operational rule not validated against surveyed, typical, or performance-tuned usage.
It is an intentionally artificial, predeclared selection policy, not a behavioral practitioner model.
The frozen pages, conflict rule, and assignments make the decision auditable, but the single author
both defined and applied it, and two blind selection reviews are prepared and still pending.
Where the policy differs from a performance-tuned configuration, RQ1's spread mixes the library's
execution machinery with that strategy choice, made material by the one-to-four statement range on
the deep fetch (Supplement Table~S2).
Three sensitivity checks qualify it, and none of them decomposes the spread: the preserved
source-page provenance is provenance, not evidence of frequency or idiomaticity; the
\emph{common-SQL raw-path sensitivity contrast} is a compound contrast, not a bound on how much
comes from the loading strategy versus the library machinery; and the loading-strategy check shows
the selected path is the \emph{faster} strategy for the join-selected ORMs, with Objection's slower
selected select-in path on MySQL disclosed rather than corrected.
RQ1 therefore answers what this policy yields, not a library-only overhead.

\paragraph{Internal validity}
Each adapter follows the registered policy and passes specification-derived and differential checks,
but all were implemented by one author.
A semantically correct yet needlessly allocating, converting, or serializing adapter would pass; no
independent human implementation review is claimed.
A result-blind author self-audit found two avoidable timed-path operations before campaign
acceptance, both removed without changing the declared treatment or SQL, after which the admission
chain passed again and measurement restarted from repetition~1
(Section~\ref{sec:discussion}); because the implementer performed it, this is author review.

Both are now detected mechanically instead, which is why R7 requires evidence rather than a reviewer:
a static check reports values acquired and never read, and a runtime gate rejects any cell writing to
\texttt{stderr} or \texttt{stdout} during a measured run.
The runtime gate postdates the reported campaign and did not gate it, so absence of output there
rests on author observation.
Neither covers waste that is silent \emph{and} invisible in the source, so the residual narrows
without closing: a correct, policy-conformant adapter uniformly slower than a differently written one
would pass every check.
We report it as open, not discharged.
The same caution applies more widely: four admission elements postdate the accepted campaigns
(Supplement Table~S48).
Because they replay a deterministic seed and inspect archived state, they can invalidate archived
results, and did so for the superseded pilot; but they did not gate the timed runs as those runs
executed, so a defect observable only while a run was in progress would have escaped them.

Per-cell warm-up phases exceeding the slowest layer's cold-start stabilization absorb JIT, pool-fill,
and plan-cache effects~\cite{mytkowicz2009producing}, and two further controls address non-layer
confounds: one correlated-subquery aggregation plan and the corrected-state campaign's verified reset
before each write cell.
The historical MySQL allocator defect, the corrected reset, and the seeded identifier design are
described in Section~\ref{sec:methodology}; their consequence here is that the affected historical
cells are not used as clean evidence, all five patterns having been remeasured in one campaign.
A forward-versus-reversed write-cell-order check found no detectable order effect, so order does not
confound layer identity with run position.
A residual risk is coordinated omission, which can under-report the tail at
saturation~\cite{fruth2022tail,tene2015latency}; the corrected constant-arrival sweeps show small
tails below measured capacity and sharp growth near it (Supplement Tables~S6 and S17).

\paragraph{External validity}
All measurements come from a single schema, dataset size, and 16-vCPU virtualized host, a pool fixed
at ten, and specific engine and library versions (PostgreSQL~18.4, MySQL~9.7.1; latest compatible
stable versions as of the 15 July 2026 freeze, Supplement Table~S11).
The fixed pool does not appear to drive the ranking: the pool-size frontier preserves the ordering
from 2 to 50 connections and a mixed read/write workload preserves the read ordering, both
exploratory and on four layers, though at a pool of one Prisma overtakes Knex on PostgreSQL
(Supplement Tables~S7, S25, S26).
Two choices are disclosed rather than resolved: Sequelize on its stable 6.x line, and Prisma's MySQL
path on the MariaDB adapter, an implementation$\times$engine asymmetry.
The memory-resident working set is a hot-cache regime that makes access-layer overhead more visible;
as a workload becomes I/O-bound the spreads may compress toward $1\times$.

Same-host checks do not bound an independent substrate: a post-restart rerun of six representative
cells moved absolute throughput materially against the campaign prevailing when it ran, which is the
superseded pilot rather than the accepted campaign.
The corrected-state 25-run campaign repeats all five patterns with a new randomized order and
environment fingerprint on the same host.
Cross-campaign agreement or disagreement is reported directly; it reduces dependence on one overnight
ordering but does not address host, hypervisor-neighbor, CPU-steal, or frequency transfer.
No empirical ordering is intended to travel beyond the tested campaigns, and close cross-backend rank
reversals are exploratory unless their direction repeats.
Every promoted RQ2 reversal involves the single treatment whose MySQL driver differs, and a
leave-one-out repetition of the promotion procedure without it promotes none (Supplement Table~S49);
the promoted-reversal evidence is therefore conditional on that driver substitution, which the design
cannot separate from the treatment or from the DBMS.
The protocol evidence is also bounded: its external audit has one rater, so general independent
applicability remains to be tested.

\paragraph{Conclusion validity}
Our analysis (Section~\ref{sec:methodology}) reports medians with the coefficient of variation and
within-campaign bootstrap 95\% intervals (repeatability, not population-general) and, because the
design is a randomized block, compares adjacently ranked layers with \emph{paired} tests on their
per-replicate throughput ratios.
Two bounds reinforce the ordering: the widest spreads dwarf the deep-fetch run-to-run variability,
while insert variability is large enough on both engines that no fine insert ordering is claimed
(Supplement Tables~S1, S9), and the ordering holds across the concurrency sweep at $\geq32$
connections (Supplement Figure~S2).
Rather than assert a specific statistical power, absent a prespecified effect-size model, we rest the
conclusions on effect sizes and interval estimates, foremost the \emph{predefined} native-driver
contrasts (Supplement Table~S41), keeping the ranking-selected adjacent-pair permutation tests only
as a descriptive post-hoc check.
Application processes do not cross adapter, engine, or replicate blocks, although read endpoints
within a block share one process~\cite{kalibera2013rigorous}.
All blocks share one host; primary intervals remain within-campaign, while the corrected RQ2 subset
has a second same-host campaign.
The secondary experiments covering only a layer subset or few replicates (Supplement Table~S32) are
read as \emph{sensitivity evidence} on the tested subset, not as claims about all eleven
implementations.
The common-SQL contrast and the matched-utilization tails rest on few replicates, so both are
reported with their observed spread and neither carries an interval estimate (Supplement
Tables~S14, S21--S22).
The matched-utilization tails carry a second, unpropagated uncertainty: each offered rate is a
fraction of an \emph{estimated} capacity, whose spread is evaluated separately (Supplement
Table~S45) rather than carried into the reported p99.
They are secondary evidence for the operating-point distinction, not a precise tail estimate.

\section{Conclusion and Future Work}
\label{sec:conclusion}

This paper proposes a concrete discipline for access-layer comparisons: derive expected read results from the declared task, check mutual equivalence and post-write state before a timing is admitted, define each configured treatment by a predeclared policy, and report fixed-demand latency separately from capacity and matched utilization. The selection policy is reproducible evidence about the treatments it fixes, not observed practitioner behavior, and the common-SQL raw-path experiment is a compound sensitivity contrast, not a library-intrinsic overhead estimate.

In the pinned Express.js experiment, the largest observed gap occurs on the tested deep fetch: $7.01\times$ on PostgreSQL and $5.12\times$ on MySQL, narrowing, in a separate ten-run campaign, to $1.71\times$ and $2.03\times$ under common SQL through raw facilities. At stable 50\% and most 70\% capacity targets, p99 values form a much smaller band than at the high-load point. At the stated evidential bar only one treatment reverses durably across the two tested backend stacks, and it is the one whose MySQL driver differs, so the experiment records limited, treatment-specific non-transfer without identifying the DBMS as its source.

The evidence record separates completed checks from unperformed validation: the oracle, preflight,
second same-host campaign, checklist, and retrospective audit make the method inspectable, while the
audit has one rater, every adapter is single-author, and both campaigns share one virtualized host.
Result-blind packets are provided instead of a claim of independent validation.

Future work should estimate inter-rater agreement for the external audit, repeat the reduced RQ2
matrix on a physically independent host, and extend the workload across datasets, topologies, and
cold-cache regimes. The first is prepared rather than proposed: a rater-extensible audit record and a
committed agreement scorer ship with the artifact.


\section*{Data availability}
The public repository is available at \texttt{https://github.com/mmiotk/express-db-access-performance}. The complete revision artifact is archived as release v1.12.15 under the immutable version DOI \texttt{https://doi.org/10.5281/zenodo.21599104}; the concept DOI is \texttt{https://doi.org/10.5281/zenodo.21313858}. The GitHub tag identifies the exact release commit. A reproducibility summary --- versions,
requirements, run commands, time and resources, and which results the harness regenerates
automatically from the archived raw data --- is in Supplement Table~S31 and
\texttt{REPRODUCE.md}. \textbf{Exact reproduction requires the reference engines}
(PostgreSQL~18.4 and MySQL~9.7.1), installed via the user-space conda path documented in
\texttt{REPRODUCE.md}; the convenience Docker path pins the same engine versions but uses relaxed durability in a containerized topology and therefore reproduces the workflow, not the headline default-durability insert numbers. A machine-readable
encoding of the protocol --- inputs, mandatory and recommended stages, the cell-admission gate, and
compliance levels --- ships as \texttt{protocol-checklist.yaml} for auditing a benchmark against the
protocol. Exact 42-file source archives preserve and verify the measurement-time
source state of each accepted campaign. Reproduction status is separated explicitly. On 26 July 2026, an author-run isolated reconstruction of the current candidate verified 60/60 archived JSON hashes and regenerated 35 table outputs and four analysis JSON files byte-for-byte without starting either DBMS. An earlier author-run archive-isolated check of immutable v1.12.9 verified 35/35 inputs and reproduced 45/50 then-committed tables; its five presentation-only differences are identified in \texttt{notes/clean-room-reproduction.md}. Neither check is independent reproduction or a re-execution of the current measurements, and neither is described as such. This is the generic build; the Elsevier submission build is
\texttt{ist/ist\_main.tex}.

\section*{CRediT authorship contribution statement}
\textbf{Mateusz Miotk:} Conceptualization, Methodology, Software, Investigation,
Data curation, Formal analysis, Visualization, Writing -- original draft,
Writing -- review \& editing.

\section*{Declaration of competing interest}
The author declares no known competing financial interests or personal
relationships that could have appeared to influence the work reported in this paper,
and is not affiliated with, nor funded by, any of the database libraries or vendors
evaluated in this study.

\section*{Funding}
This research did not receive any specific grant from funding agencies in the
public, commercial, or not-for-profit sectors.

\section*{Declaration of generative AI and AI-assisted technologies in the writing process}
During the preparation of this work the author used large language models (Claude, Anthropic, and Codex, OpenAI) in order to draft and edit the manuscript text and improve its readability and
language. After using these tools, the author reviewed and edited the content as needed and
takes full responsibility for the content of the published article.

\section*{AI-assisted research software}
Separately from the writing above, large language models (Claude, Anthropic, and Codex, OpenAI) also assisted
the author in implementing and analyzing the benchmark harness. This is research tooling for
which the author takes full responsibility, and the AI-assisted analytical code is publicly inspectable and checked: the statistical estimators carry unit tests against
hand-computed values and known closed-form properties (\texttt{bench/stats.test.mjs}), every
reported table cell traces to a raw per-cell record through a committed generator
(\texttt{MANIFEST.md}), and the specification-derived read oracle checks expected results without using timed native-driver output; differential checks then establish finite cross-implementation equivalence; and write checks validate resulting database state before timing is trusted.

\printbibliography

\end{document}
